\section{La transformación de XPeng (小鹏) hacia Physical AI: ciclo de datos del fabricante de automóviles y objetivo de Robotaxi}

Las secciones anteriores trataron la pila tecnológica y la eliminación de Language; esta sección examina por qué XPeng elevó la conducción autónoma al nivel de Physical AI. Liu Xianming (刘先明) afirma que «XPeng es, en esencia, una empresa de IA». Esta frase no es un eslogan publicitario; significa lo siguiente: el automóvil es uno de los agentes inteligentes del mundo físico de mayor escala, más controlables y con un ciclo comercial más completo; el fabricante de automóviles cuenta con usuarios reales, carreteras reales, retroalimentación real y una cadena de hardware controlable.

\begin{figure}[H]
\centering
\includegraphics[width=0.90\textwidth]{physical-ai-strategy.png}
\caption{La transformación de XPeng hacia Physical AI: de la capacidad de conducción autónoma a la inteligencia en el mundo físico. Diagrama conceptual propio, elaborado a partir de la conversación de 00:33:53--00:54:30.}
\end{figure}

\begin{knowledgebox}{Lectura de la figura: el automóvil es uno de los mayores cuerpos físicos de Physical AI}
En la figura, Physical AI no comprende solo el modelo, sino también el automóvil, los datos, el Robotaxi, la organización y la producción en serie. La ventaja de XPeng no está únicamente en los algoritmos, sino en que, como fabricante de automóviles, puede integrar los datos reales de su flota y la retroalimentación de la producción en serie en la iteración del modelo.
\end{knowledgebox}

\subsection{La ventaja en datos del fabricante de automóviles}

La mayor ventaja del fabricante de automóviles es una cadena de datos controlable. La flota circula en ciudades reales y genera datos de sensores, decisiones de conducción, retroalimentación de los usuarios, corner case anómalos, tomas de control y órdenes de incidencias. Siempre que el retorno de datos y la infraestructura de entrenamiento sean lo bastante buenos, el fabricante puede formar el ciclo «datos del mundo real -> entrenamiento en la nube -> despliegue en el vehículo -> retroalimentación real».

\begin{figure}[H]
\centering
\includegraphics[width=0.96\textwidth]{data-infra-loop-xpeng.png}
\caption{Ciclo de datos e Infra: la ventaja del fabricante de automóviles está en la cadena de datos controlable y en la retroalimentación de la producción en serie. Diagrama conceptual propio, elaborado a partir de la conversación de 00:02:16--00:19:00 y 00:33:53--00:54:30.}
\end{figure}

\begin{knowledgebox}{Lectura de la figura: el ciclo de datos no se reduce a «tener automóviles»}
La flota es solo el punto de entrada de los datos; un ciclo real requiere además la detección de corner case, infra de entrenamiento y análisis, actualización del modelo, despliegue en el vehículo y validación en producción en serie. Sin este sistema, los datos no se convierten por sí solos en capacidad.
\end{knowledgebox}

\subsection{El Robotaxi como hito de etapa}

Liu Xianming mencionó que espera que, en uno a tres años, el Robotaxi funcione muy bien en Guangzhou. El sentido de este objetivo es convertir el debate técnico en un servicio real: deja de ser una noticia del momento, un juguete o una atracción turística, y pasa a formar parte de lo que los usuarios usan por defecto cada día para desplazarse. La transición de Waymo en San Francisco, de atracción turística a servicio cotidiano, es la referencia que más le importa.

\begin{knowledgebox}{De atracción turística a servicio cotidiano}
Cuando Liu Xianming viajaba en Waymo con sus hijos en San Francisco, los niños lo llamaban Robot Car. Al principio, mucha gente iba a San Francisco solo para probarlo por novedad; pero cuando empezó a atender una gran cantidad de viajes diarios, y al salir de Caltrain Station ya había personas esperando un vehículo, dejó de ser solo una demostración tecnológica y pasó a formar parte de la vida de la ciudad. El objetivo de Robotaxi de XPeng también tiene que cruzar ese umbral de percepción.
\end{knowledgebox}

\subsection{Por qué Physical AI no es solo conducción autónoma}

Esta sección sitúa el objetivo de Robotaxi dentro del marco más amplio de Physical AI. La conducción autónoma es uno de los escenarios de la IA física más maduros, con más datos y con el ciclo comercial más claro, pero no es el destino final. El sistema del vehículo ya incluye percepción, predicción, planificación, control, interacción con el usuario, mapas, entrenamiento en la nube y mecanismos de seguridad; en el futuro, estos módulos se trasladarán a más agentes inteligentes físicos. Si XPeng logra demostrar el scaling en el cuerpo físico del automóvil, tendrá la oportunidad de extender la misma metodología a tareas más amplias del mundo físico.

\begin{importantbox}{El automóvil es el campo de entrenamiento de Physical AI}
El automóvil es un agente inteligente de alto valor, alta frecuencia y fuertes restricciones del mundo real. Permite que la IA no solo responda preguntas en una pantalla, sino que aprenda a actuar en medio de carreteras, tráfico y responsabilidades de seguridad.
\end{importantbox}

\subsection{Resumen de la sección}

Lo esencial de la transformación de XPeng hacia Physical AI es situar la conducción autónoma dentro de un ciclo de inteligencia del mundo real: los datos del fabricante de automóviles, el entrenamiento en la nube, el despliegue en el vehículo, la validación con Robotaxi y la ejecución organizacional conforman juntos la nueva estrategia.
